\documentclass[11pt,a4paper]{article}

% ─── Encodage & langue ───────────────────────────────────────────────────────
\usepackage[utf8]{inputenc}
\usepackage[T1]{fontenc}
\usepackage[french]{babel}
\usepackage{lmodern}

% ─── Mise en page ────────────────────────────────────────────────────────────
\usepackage[margin=2.5cm]{geometry}
\usepackage{microtype}
\usepackage{graphicx}
\usepackage{booktabs}
\usepackage{array}
\usepackage{enumitem}
\usepackage{amsmath}
\usepackage{xcolor}
\usepackage{listings}
\usepackage{hyperref}
\hypersetup{colorlinks=true, linkcolor=black, urlcolor=blue, citecolor=black}

% ─── Style code ──────────────────────────────────────────────────────────────
\definecolor{codebg}{rgb}{0.96,0.96,0.96}
\definecolor{codekw}{rgb}{0.10,0.30,0.65}
\definecolor{codecm}{rgb}{0.30,0.55,0.30}
\lstset{
  basicstyle=\ttfamily\footnotesize,
  backgroundcolor=\color{codebg},
  keywordstyle=\color{codekw}\bfseries,
  commentstyle=\color{codecm}\itshape,
  stringstyle=\color{purple},
  breaklines=true, frame=single, framesep=4pt,
  showstringspaces=false, columns=fullflexible,
  numbers=left, numberstyle=\tiny\color{gray}, numbersep=6pt,
}

\newcommand{\code}[1]{\texttt{\small #1}}

% ─── Page de titre ───────────────────────────────────────────────────────────
\title{\textbf{Conception et développement d'un système de suivi\\
d'activité utilisateur pour visualisations web}\\[4pt]
\large Composant de suivi du regard de l'\textit{Activity Tracker}}
\author{Projet PIR — Laboratoire i3S, Université Côte d'Azur\\
\small Encadrante : Aline Menin}
\date{Juin 2026 \\ \small\textit{Version~1 — rapport technique (brouillon)}}

\begin{document}
\maketitle

\begin{abstract}
\noindent
Ce rapport décrit le prototype du composant de \textbf{suivi du regard}
(\textit{eye-tracking}) d'un \textit{Activity Tracker} destiné à l'étude des
processus de raisonnement analytique sur des visualisations de données web. Le
prototype repose sur \textbf{WebGazer.js} pour estimer la direction du regard à
partir d'une webcam standard, sans matériel spécialisé. Il fournit~: un module de
capture, une calibration avancée (correction spatiale continue, compensation des
mouvements de tête, validation par leave-one-out), un post-traitement extrayant
\textbf{fixations} et \textbf{saccades} (algorithmes I-DT et I-VT), et un
\textbf{format de journalisation structuré} exporté en JSON, JSON-LD et CSV pour
intégration ultérieure dans un graphe de connaissances. Deux moteurs de regard
interchangeables sont proposés au choix de l'utilisateur (WebGazer.js et un moteur
maison fondé sur MediaPipe FaceLandmarker). La validation logicielle
s'appuie sur 506 tests unitaires automatisés. Ce document constitue une première
version~: la section d'évaluation utilisateur (3--5 participants) présente le
protocole et un cadre d'analyse à compléter par les données réelles.
\end{abstract}

\tableofcontents
\bigskip\hrule

% ═════════════════════════════════════════════════════════════════════════════
\section{Introduction et contexte}

Comprendre comment les utilisateurs explorent visuellement les visualisations de
données est essentiel pour étudier le raisonnement analytique. En recherche en
visualisation d'information, les journaux d'interaction et le suivi du regard sont
couramment combinés pour analyser le comportement et l'attention visuelle.

Les progrès du suivi du regard \emph{dans le navigateur}, en particulier
WebGazer.js~\cite{papoutsaki2016}, permettent d'estimer le regard depuis une
webcam standard, ce qui facilite les études en ligne à grande échelle. Ce projet
s'inscrit dans un effort plus large visant à construire un \textit{Activity
Tracker} collectant des données multimodales (interactions, regard,
verbalisations) pour étudier les processus analytiques et soutenir la conception
de visualisations adaptatives, les données étant à terme intégrées dans un
\textbf{graphe de connaissances}.

\paragraph{Objectif.} Prototyper le composant de suivi du regard dans un
environnement de visualisation web. Les tâches sont~: (T1) étudier les solutions
open-source~; (T2) implémenter la capture~; (T3) extraire fixations et saccades~;
(T4) définir un format de journalisation~; (T5) tester avec 3--5 participants.

% ═════════════════════════════════════════════════════════════════════════════
\section{T1 — Étude des solutions open-source}

Plusieurs bibliothèques de suivi du regard par webcam existent. Le tableau~\ref{tab:compare}
résume les principales options évaluées.

\begin{table}[h]
\centering
\renewcommand{\arraystretch}{1.25}
\begin{tabular}{>{\raggedright\arraybackslash}p{2.7cm} p{3.0cm} p{3.0cm} p{3.0cm}}
\toprule
\textbf{Critère} & \textbf{WebGazer.js} & \textbf{MediaPipe FaceLandmarker} & \textbf{GazeCloudAPI} \\
\midrule
Exécution & 100\,\% navigateur & 100\,\% navigateur (WASM/GPU) & Service cloud \\
Matériel & Webcam standard & Webcam standard & Webcam standard \\
Licence & Open-source (GPLv3) & Open-source (Apache) & Propriétaire \\
Vie privée & Local & Local & Vidéo au cloud \\
Sortie & Point écran direct & Landmarks + pose (mapping à coder) & Point écran \\
Calibration & Intégrée & À implémenter (faite ici) & Intégrée \\
Maturité & Élevée, citée en recherche & Modèle récent, précis & Moyenne \\
\bottomrule
\end{tabular}
\caption{Comparaison des solutions de suivi du regard par webcam.}
\label{tab:compare}
\end{table}

\paragraph{Choix : les deux moteurs locaux.} GazeCloudAPI est écarté (vie privée~:
il transmet la vidéo à un tiers). Deux moteurs \emph{locaux} sont retenus et
proposés au choix de l'utilisateur via une page d'accueil~:
\begin{itemize}[leftmargin=1.4em]
  \item \textbf{WebGazer.js} — solution clé-en-main, mapping regard$\to$écran intégré,
  reproductible et citée en recherche~; idéal pour un prototype rapide.
  \item \textbf{MediaPipe FaceLandmarker} — détection d'iris et pose de tête plus
  précises, mais ne fournit pas de point de regard~: tout le mapping a été
  développé (cf. §\ref{sec:mediapipe}).
\end{itemize}
Proposer les deux permet de comparer empiriquement précision et robustesse sur le
même protocole et le même format de données.

% ═════════════════════════════════════════════════════════════════════════════
\section{T2 — Module de capture du regard}

\subsection{Architecture}

Le prototype est une application web statique (HTML + JavaScript, sans étape de
\textit{build}), organisée en modules indépendants exposés comme objets globaux
(tableau~\ref{tab:modules}).

\begin{table}[h]
\centering
\begin{tabular}{l l p{7.5cm}}
\toprule
\textbf{Fichier} & \textbf{Global} & \textbf{Rôle} \\
\midrule
\code{gaze-capture.js} & \code{GazeCapture} & Démarrage/arrêt de WebGazer, émission des points. \\
\code{calibration.js}  & \code{Calibration} & Calibration, correction spatiale, lissage, détection fixations/saccades. \\
\code{gaze-logger.js}  & \code{GazeLogger}  & Journalisation et export structuré (JSON / JSON-LD). \\
\bottomrule
\end{tabular}
\caption{Modules principaux et API publiques.}
\label{tab:modules}
\end{table}

\subsection{Capture}

Le module \code{GazeCapture} vérifie la disponibilité de la webcam via
\code{getUserMedia}, initialise WebGazer (régression \emph{ridge}, traqueur
TFFacemesh), puis émet à chaque trame un point $\{x, y, \text{timestamp}\}$ aux
abonnés. Deux précautions améliorent la qualité du modèle~: la suppression des
écouteurs souris (\code{removeMouseEventListeners}), car le mouvement de souris
pollue l'entraînement, et la contrainte explicite de résolution caméra.

% ═════════════════════════════════════════════════════════════════════════════
\section{Calibration et précision}

Au-delà du sujet, un effort particulier a porté sur la \textbf{qualité de
calibration}, principal facteur limitant du suivi du regard par webcam.

\subsection{Procédure}

Une grille de \textbf{25 points} ($5\times5$, bords densifiés) est présentée~; à
chaque point, plusieurs clics enregistrent des couples (regard, position écran).
Une grille de \textbf{9 points indépendants} sert à la validation. Un écran de
positionnement préalable vérifie luminosité, équilibre d'éclairage et distance
($\approx$60\,cm) via l'analyse de la luminance webcam.

\subsection{Correction spatiale}

L'erreur résiduelle après régression n'est pas uniforme sur l'écran. Un
\textbf{champ de correction continu} est estimé à partir des points de validation.
Une régression locale pondérée \textbf{LOWESS} (noyau tricubique) corrige chaque
prédiction, avec repli sur une pondération par distance inverse (IDW) hors de la
bande passante~:
\[
\hat{e}(p) = \frac{\sum_i w_i\, e_i}{\sum_i w_i}, \qquad
w_i = \left(1 - \left(\tfrac{d_i}{h}\right)^3\right)^3 \ \text{pour}\ d_i < h .
\]
Un point d'attention important a été corrigé~: une version antérieure
\emph{chaînait} IDW puis LOWESS, tous deux estimant le \emph{même} champ d'erreur,
ce qui produisait une \textbf{sur-correction}. Un seul estimateur est désormais
appliqué.

\subsection{Score honnête par validation croisée}

Évaluer la précision sur les points qui ont servi à construire la correction est
optimiste (évaluation sur le jeu d'entraînement). Le prototype calcule donc une
erreur de \textbf{généralisation par leave-one-out}~: pour chaque point de
validation, le champ de correction est reconstruit \emph{sans} ce point, puis
l'erreur résiduelle est mesurée sur le point exclu. C'est cette valeur qui est
rapportée comme indicateur réaliste.

\subsection{Lissage temporel et compensation de tête}

En temps réel, un filtre \textbf{One Euro}~\cite{casiez2012} lisse le tremblement
au repos sans introduire de latence sur les saccades. WebGazer supposant la tête
immobile, une \textbf{compensation des mouvements de tête} a été ajoutée~: le
centroïde des points caractéristiques du visage est mémorisé en référence, et tout
déplacement ultérieur génère un décalage écran proportionnel et borné, réduisant la
dérive sur les sessions longues.

% ═════════════════════════════════════════════════════════════════════════════
\section{Second moteur : MediaPipe FaceLandmarker}
\label{sec:mediapipe}

Un second moteur, alternatif et interchangeable avec WebGazer, a été développé sur
\textbf{MediaPipe FaceLandmarker}. Contrairement à WebGazer, MediaPipe ne fournit
pas un point de regard à l'écran mais 478 points caractéristiques 3D (iris inclus)
et une matrice de pose de tête. Toute la chaîne de régression
\emph{caractéristiques$\to$écran} a donc été construite~:

\begin{enumerate}[leftmargin=1.5em]
  \item \textbf{Extraction de caractéristiques}~: position relative de l'iris dans
  chaque orbite, exprimée dans un \emph{repère frontal canonique} obtenu en annulant
  la rotation de tête (projection par $R^{\top}$)~; angles d'Euler de la pose
  (yaw/pitch/roll)~; échelle inter-oculaire (proxy de distance)~; termes polynomiaux
  et blendshapes oculaires.
  \item \textbf{Régression ridge}~: apprentissage par moindres carrés régularisés sur
  caractéristiques \emph{standardisées}, avec pondération des échantillons par leur
  qualité et choix automatique du paramètre $\lambda$ par validation croisée k-fold.
  \item \textbf{Correction résiduelle et lissage}~: champ de correction
  (IDW/LOWESS) appris sur la validation, filtre One~Euro temps réel, et
  \emph{apprentissage continu} (recalibration incrémentale par clics en cours
  d'usage).
\end{enumerate}

Ce moteur offre une compensation des mouvements de tête \emph{native} (la pose fait
partie des caractéristiques) et une cadence élevée grâce à l'inférence GPU
synchronisée sur les images de la webcam. Les deux moteurs partagent la même chaîne
de post-traitement, le même journal et le même format d'export, ce qui rend leur
comparaison directe sur le même protocole.

\paragraph{Deux modes de calibration.} Outre la calibration classique par clics
(grille de 25 points), un mode \textbf{calibration par poursuite} (\emph{smooth
pursuit}) a été ajouté pour MediaPipe~: une cible mobile parcourt l'écran en
serpentin et l'utilisateur la suit du regard, sans cliquer. Ce mode collecte
beaucoup plus d'échantillons (plusieurs centaines) le long d'une trajectoire
continue, ce qui produit souvent un modèle de régression plus stable.

% ═════════════════════════════════════════════════════════════════════════════
\section{T3 — Extraction des fixations et saccades}

Deux algorithmes classiques de l'analyse du regard~\cite{salvucci2000} sont
implémentés.

\paragraph{Fixations — I-DT (Dispersion-Threshold).} Une fenêtre temporelle
glissante regroupe les points dont la dispersion spatiale
$(\max_x - \min_x) + (\max_y - \min_y)$ reste sous un seuil pendant une durée
minimale. L'implémentation utilise des \textit{deques} monotones pour calculer la
dispersion en $O(n)$ amorti~; le traitement de 10\,000 points s'exécute en moins
de 50\,ms.

\paragraph{Saccades — I-VT (Velocity-Threshold).} Les transitions dont la vélocité
instantanée dépasse un seuil sont marquées comme saccades, caractérisées par leur
amplitude et leur vélocité maximale. Une fonction \code{linkEvents} produit une
chronologie alternant fixations et saccades.

Chaque fixation peut être associée à une \textbf{zone d'intérêt} (AOI) de la
visualisation regardée, base de l'analyse attentionnelle.

% ═════════════════════════════════════════════════════════════════════════════
\section{T4 — Format de journalisation structuré}

\subsection{Structure}

Une session exportée comporte cinq sections (listing~\ref{lst:format})~: métadonnées
\code{session}, points bruts \code{raw\_gaze\_data}, événements \code{events}
(fixations/saccades), \code{aoi\_hits}, et \code{interactions} multimodales. Le
format est décrit par un \textbf{JSON Schema} (draft-07) qui le rend validable et
auto-documenté.

\begin{lstlisting}[language=, caption={Structure d'une session exportée (extrait).}, label=lst:format]
{
  "session": {
    "id": "uuid", "format_version": "1.1.0",
    "participant_id": "P01", "clock_origin_ms": 12345.6,
    "calibration_score": { "mean_error_px": 92.4 },
    "config_snapshot": { "lowess_bandwidth": 0.45, ... }
  },
  "raw_gaze_data": [ { "x": 812, "y": 437, "t_rel_ms": 4502.1 }, ... ],
  "events":        [ { "type": "fixation", "duration": 280,
                       "details": { "x": 810, "y": 440 } }, ... ],
  "aoi_hits":      [ { "aoi_id": "bar-q3", "event_index": 12 }, ... ],
  "interactions":  [ { "type": "tab_change", "t_rel_ms": 8100.0 }, ... ]
}
\end{lstlisting}

\subsection{Choix de conception}

\begin{itemize}[leftmargin=1.4em]
  \item \textbf{Double horloge.} Chaque enregistrement porte un \code{timestamp}
  epoch (lisible) \emph{et} un \code{t\_rel\_ms} issu de \code{performance.now()},
  horloge \textbf{monotone} indispensable au calcul fiable des durées et vélocités.
  \item \textbf{Reproductibilité.} \code{config\_snapshot} fige les paramètres de
  filtrage/correction effectifs, permettant de rejouer ou comparer deux sessions.
  \item \textbf{Multimodalité.} La section \code{interactions} journalise clics,
  changements d'onglet et entrées du regard dans une AOI, conformément à l'objectif
  d'\textit{Activity Tracker}.
\end{itemize}

\subsection{Richesse contextuelle de chaque point de regard}
\label{sec:richesse}

Au-delà des coordonnées, chaque point de regard est enrichi de tout le contexte
permettant de décrire \emph{ce que l'utilisateur regardait}~:

\begin{itemize}[leftmargin=1.4em]
  \item \textbf{Confiance} de la prédiction ($\in[0,1]$) et \textbf{coordonnées
  brutes} avant correction, en plus des coordonnées corrigées.
  \item \textbf{Module source} (\code{webgazer} ou \code{mediapipe}) sur chaque
  entrée, rendant chaque module clairement identifiable dans le journal.
  \item \textbf{Descripteur DOM} de l'objet observé~: balise, identifiant, classes,
  \emph{type sémantique} (barre, axe, point, légende, courbe, étiquette), boîte
  englobante, texte et attributs \code{data-*}, sélecteur CSS. On sait ainsi si le
  regard portait sur une barre, un axe, etc.
  \item \textbf{État de la visualisation}~: vue active, jeu de données, niveau de
  zoom, filtres actifs, sélection et AOI courante. Un historique \code{viz\_states}
  enregistre en plus chaque changement d'état (zoom, filtre par trimestre, onglet).
\end{itemize}

Chaque point porte en outre la \textbf{luminosité ambiante} (\code{lux}, mesurée par
l'API \textit{AmbientLightSensor} ou estimée depuis la webcam) et, pendant un test
guidé, le \textbf{contexte du stimulus} (\code{test\_case\_id}, \code{target\_aoi\_id}).
La session enregistre les informations participant (\code{first\_name},
\code{last\_name}, \code{glasses}, \code{engine}). Le format est versionné (v1.3).

Trois formats d'export sont fournis~: \textbf{JSON} (structuré), \textbf{JSON-LD}
(graphe de connaissances) et \textbf{CSV} (analyse tabulaire directe sous
pandas/R/Excel).

\subsection{Export orienté graphe de connaissances (JSON-LD)}

Un second export en \textbf{JSON-LD} représente la session comme un graphe de
nœuds typés sous un vocabulaire \code{wga:} (\code{wga:Session}, \code{wga:Fixation},
\code{wga:Saccade}, \code{wga:AOIHit}, \code{wga:Interaction}). Ce format se charge
directement dans un triple store RDF, anticipant l'intégration au graphe de
connaissances de l'\textit{Activity Tracker}.

% ═════════════════════════════════════════════════════════════════════════════
\section{T5 — Évaluation avec participants}

\subsection{Protocole — session de test guidée et automatique}

Les deux pages moteurs (\code{index-webgazer.html} et \code{index-mediapipe.html})
implémentent le même protocole \textbf{entièrement guidé et automatique}, sans aucune
action de l'utilisateur pendant le test~:
\begin{enumerate}[leftmargin=1.5em]
  \item \textbf{Formulaire participant}~: prénom, nom, port de lunettes.
  \item \textbf{Calibration} (par clics ou, pour MediaPipe, par poursuite) puis
  lancement \emph{automatique} du scénario.
  \item \textbf{Scénario}~: enchaînement automatique des trois visualisations
  (page~1 bar chart, page~2 line chart, page~3 scatter). Sur chaque page,
  \textbf{trois cercles rouges} « regardez ici » apparaissent successivement
  ($\approx$3\,s chacun) sur des AOI réelles distinctes, puis une phase
  d'\textbf{exploration libre} de 15\,s, avant le passage automatique à la page
  suivante. Chaque point est étiqueté avec le stimulus, l'AOI cible et sa position
  écran (\code{test\_case\_id}, \code{target\_aoi\_id}, \code{target\_x/y}).
  \item \textbf{Résultats}~: post-traitement, enregistrement de la session et
  affichage de l'écran de résultats détaillés (analyse + parcours).
\end{enumerate}
Un \textbf{retour visage permanent} (visage détecté, distance, luminosité) reste
affiché en continu, y compris pendant le test, pour signaler tout problème de
captation. Chaque session produit \textbf{un fichier JSON par personne} et est
persistée localement pour reconsultation.

\paragraph{Mesures collectées.} Erreur de calibration (px), luminosité (lux),
confiance moyenne, nombre de fixations/saccades, hits AOI, et — grâce à
l'étiquetage des stimuli — la correspondance entre l'AOI ciblée et l'objet
réellement regardé.

\paragraph{À compléter.} Les sessions réelles (3--5 participants) seront collectées
via cette page et consultées dans le visualiseur (§\ref{sec:viewer})~: tableau
récapitulatif par participant (nom, date, lunettes, lux, erreur de calibration),
distribution des erreurs et synthèse qualitative.

% ═════════════════════════════════════════════════════════════════════════════
\section{Visualisation et consultation des sessions}
\label{sec:viewer}

Une sous-page \code{session-viewer.html} permet de \textbf{consulter et comparer}
les sessions enregistrées sans relancer de test. Les sessions étant persistées
localement (le projet étant statique, \code{localStorage} sert de base de données ;
un export/import fichier permet l'archivage et le transfert), la page offre~:
\begin{itemize}[leftmargin=1.4em]
  \item un \textbf{tableau triable et filtrable} (par moteur, par port de lunettes)
  de toutes les sessions (date, participant, moteur, erreur de calibration, lux,
  confiance, nombre de points/fixations/saccades)~;
  \item un \textbf{écran de résultats détaillés} en deux onglets : une \emph{analyse
  par stimulus} (page regardée, AOI ciblée, score « sur la cible », durée) et un
  \emph{parcours du regard} montrant les \textbf{trois visualisations re-dessinées}
  avec, par-dessus chacune, la \textbf{heatmap} de densité (vert$\to$orange$\to$rouge)
  et le scanpath des fixations propres à cette page~;
  \item une \textbf{comparaison} entre deux sessions ou deux groupes (p.\,ex.\ toutes
  les sessions WebGazer contre toutes les MediaPipe), présentée sous forme d'un
  \textbf{bilan écrit en français} (quel moteur est meilleur et pourquoi) suivi d'un
  tableau détaillé métrique par métrique (valeur, écart, gagnant)~;
  \item l'\textbf{export} par session (JSON, CSV) et global.
\end{itemize}

\paragraph{Score « sur la cible » continu.} Pendant le test, la position écran du
stimulus (cercle rouge) est journalisée avec chaque point de regard. Le score d'une
cible est la moyenne, sur ses points, d'une fonction de proximité valant 100\,\% si
le regard est dans un rayon serré autour du centre et décroissant continûment avec
la distance — bien plus informatif qu'un simple « dans/hors de l'AOI ».

% ═════════════════════════════════════════════════════════════════════════════
\section{Validation logicielle}

La logique métier (algorithmes I-DT/I-VT, filtres, correction spatiale, format de
log) est couverte par \textbf{506 tests unitaires} automatisés (Node, sans webcam),
exécutables par \code{npm test}. Ils vérifient notamment~: la correction unique
(non sur-corrigée), la validation leave-one-out, le lissage One Euro, la robustesse
des détecteurs sur cas limites, et la performance ($<$50\,ms sur 10\,000 points).

% ═════════════════════════════════════════════════════════════════════════════
\section{Limitations et perspectives}

\paragraph{Limitations.} La précision reste dépendante de l'éclairage, de la qualité
de la webcam et de l'immobilité de la tête — limites intrinsèques au suivi du
regard par webcam. La fréquence d'échantillonnage varie selon la machine
(typiquement 10--30\,Hz). Une dérive temporelle subsiste sur les longues sessions,
seulement atténuée par la compensation de tête.

\paragraph{Perspectives.}
\begin{itemize}[leftmargin=1.4em]
  \item Réaliser les sessions réelles (T5) et consolider l'analyse quantitative.
  \item Étendre la multimodalité aux \textbf{verbalisations} (think-aloud) prévues
  par l'\textit{Activity Tracker}.
  \item Définir et publier le vocabulaire \code{wga:} et charger les exports JSON-LD
  dans le graphe de connaissances cible.
  \item Étalonner empiriquement le gain de compensation de tête sur plusieurs
  configurations matérielles.
\end{itemize}

% ═════════════════════════════════════════════════════════════════════════════
\section{Conclusion}

Le prototype répond aux tâches T1--T4~: une solution open-source a été choisie et
justifiée, la capture, le post-traitement (fixations/saccades) et un format de
journalisation structuré (JSON et JSON-LD) sont opérationnels et testés. La
calibration a fait l'objet d'améliorations dépassant le cadre initial (correction
spatiale continue, score honnête par leave-one-out, compensation de tête). Il reste
à mener l'évaluation utilisateur (T5) pour quantifier en conditions réelles la
stabilité de la capture et les limites pratiques de l'approche.

% ═════════════════════════════════════════════════════════════════════════════
\begin{thebibliography}{9}
\bibitem{papoutsaki2016}
A.~Papoutsaki \emph{et al.},
\textit{WebGazer: Scalable Webcam Eye Tracking Using User Interactions},
IJCAI, 2016.

\bibitem{casiez2012}
G.~Casiez, N.~Roussel, D.~Vogel,
\textit{1\,€ Filter: A Simple Speed-based Low-pass Filter for Noisy Input in
Interactive Systems}, ACM CHI, 2012.

\bibitem{salvucci2000}
D.~D.~Salvucci, J.~H.~Goldberg,
\textit{Identifying Fixations and Saccades in Eye-Tracking Protocols},
ETRA, 2000.
\end{thebibliography}

\end{document}
